\chapter{A schedule that runs}\label{ch:a-schedule-that-runs}

The shortest Test Schedule that does something is eight lines. Every example in this manual
is a real file in \texttt{docs/manuals/examples/test-writer/}, and every one of them runs
headless in the test suite (\texttt{tests/test\_manuals.py}). Read this one, then run it.

\noindent\textbf{\texttt{first.tslua}}\par
\lstinputlisting[language=ATELua]{../examples/test-writer/first.tslua}

Three things are happening here.

\begin{itemize}
  \item \texttt{ScheduleParams(name, description)} builds the schedule. The name is required;
    it is what the Operator sees on the Run screen.
  \item The \texttt{setSchedule} call gives the schedule its Test Steps, in the order they
    run, and one \texttt{StepParams(name)} is one step.
  \item \texttt{:setAction(function() \dots end)} is the step's logic. ATE calls it; inside it
    the schedule calls the \texttt{ATE} table.
\end{itemize}

The setters all return the object they were given, so they chain. Order does not matter: the
setters only stash values, and ATE checks the whole schedule when the file finishes loading.

\section{Running it}

Point a Station at the examples directory, then run it from the GUI or headless. Headless, from
a source checkout, there is one script beside this manual that runs any Test Schedule any
Media Release contains:

\begin{lstlisting}[language=ATEShell]
QT_QPA_PLATFORM=offscreen python docs/manuals/examples/test-writer/run_your_schedule.py \
  --release docs/manuals/examples/test-writer --schedule "Manual: first schedule"
\end{lstlisting}

\noindent\textbf{\texttt{run\_your\_schedule.py}}\par
\lstinputlisting[language=ATEPython, caption={\texttt{docs/manuals/examples/test-writer/run\_your\_schedule.py}}]{../examples/test-writer/run_your_schedule.py}

A Step Input has to be answered or the run never ends, so each one is given with
\texttt{--input}, keyed by the \textbf{Step Name}: \texttt{--input "Operator input=1.5"}. A
Stop-on-Fail pause is answered with \texttt{--choice CONTINUE}, \texttt{RETRY} or
\texttt{ABORT}, in order. A POLL input is clicked with \texttt{--poll "Step name=BUTTON"}, which
clicks that step's first POLL input; a step with more than one POLL input names the one it means,
\texttt{--poll "Step name=2=STOP"}.
The script prints the run row, every step's Result, the Actions, the Equipment count, the
WARNING and ERROR log lines, and the Test Report paths. An answer the step's own declaration
rejects is not recorded: the script exits 1 naming the value and the rule it broke.

The shipped example release can also be run whole:

\begin{lstlisting}[language=ATEShell]
QT_QPA_PLATFORM=offscreen python examples/harness/run_examples.py
\end{lstlisting}

That harness runs the shipped example release with its simulator peers; the manual's own
examples are run by \texttt{pytest tests/test\_manuals.py}. Either way, a run writes an HTML
Test Report under \texttt{reports.directory}, and a PDF as well unless the Station is running
elevated and the Operator declines the sandbox question.

An installed ATE has no headless runner. \texttt{ate \_selftest --run} runs a built-in one-step
schedule to prove a bundle can run a Test Run, and \texttt{ate \_selftest --scan DIR} loads
every schedule in a directory and reports the broken ones --- a load check, not a run. To run
your own schedule from a command line you need a source checkout.

Two traps this script makes visible. An unanswered Step Input does not fail the schedule: the
worker blocks, and the harness raises \texttt{TimeoutError: schedule never finished}, leaving
the run row open with no \texttt{ended\_at} for the next launch to close as
\texttt{INTERRUPTED}/\texttt{FAIL}. And a POLL is answered by a click, not by a value:
\texttt{--poll}, never \texttt{--input}.

\section{What the run record shows}

A Test Run is one execution of one schedule against one UUT at one moment. ATE records the run
itself, every step, every Attempt of every step, every Operator input, every Action, every
Equipment declaration and every log line. The Test Report is generated from those rows, so
anything the schedule does that is worth knowing has to reach one of them: a
\texttt{ATE.report} call, an \texttt{ATE.log} line, or an \texttt{ATE.addEquipment} call.

\cando{write, load and run a Test Schedule of your own, and say which rows of the Test Run
record each part of it produced.}
